% !TeX root = part7.tex
\documentclass[../final.tex]{subfiles}

\begin{document}

% Продолжение по плану семинаров: тема 5; сборник, задача 4.9.
\practice{7}{17.10.2026}{Теплоёмкость кристаллической решётки}
\captionsetup{font=footnotesize,labelfont=bf,skip=5pt}

\rtheory{Энергия фононов и вклад одной моды}

В гармоническом приближении нормальные моды независимы:
\[
    \hat H=\sum_{\boldsymbol k,s}\hbar\omega_{\boldsymbol k s}
       \left(\hat n_{\boldsymbol k s}+\frac12\right).
\]
Обозначим $x_{\boldsymbol k s}=\hbar\omega_{\boldsymbol k s}/(k_{\mathrm B}T)$.
При фиксированном объёме частоты в этой модели от температуры не зависят,
и потому нулевая энергия в теплоёмкость не входит:
\begin{align*}
    U&=U_0+\sum_{\boldsymbol k,s}
       \frac{\hbar\omega_{\boldsymbol k s}}{e^{x_{\boldsymbol k s}}-1},\\
    C_V&=\left(\frac{\partial U}{\partial T}\right)_V
       =k_{\mathrm B}\sum_{\boldsymbol k,s}
          \frac{x_{\boldsymbol k s}^2e^{x_{\boldsymbol k s}}}
               {(e^{x_{\boldsymbol k s}}-1)^2}.
\end{align*}
Для одной моды при $x\ll1$ получаем $C\to k_{\mathrm B}$,
а при $x\gg1$ --- $C\simeq k_{\mathrm B}x^2e^{-x}$.
Вклад дают преимущественно моды с энергиями порядка $k_{\mathrm B}T$
или меньше.

\begingroup
\renewcommand{\thesection}{4.\arabic{section}}
\task[9]{Акустические фононы в одном, двух и трёх измерениях}
\label{task:fksv_debye}

\condition
Найти фононную энергию и теплоёмкость гармонической решётки.
Рассмотреть $d=1,2,3$, низкие и высокие температуры.
В низкотемпературном пределе использовать акустические ветви
с дисперсиями $\omega_s=c_sk$.

\solution

\subsection*{Плотность состояний}

Для системы длины, площади или объёма $V_d$ число волновых векторов
в сферическом слое $[k,k+dk]$ на одну поляризацию равно
\[
    d\mathcal N_s=\frac{V_d S_{d-1}}{(2\pi)^d}k^{d-1}dk.
\]
При $k=\omega/c_s$ получаем
\[
    D_s(\omega)=\frac{d\mathcal N_s}{d\omega}
       =\frac{V_d S_{d-1}}{(2\pi)^d}
          \frac{\omega^{d-1}}{c_s^d}.
\]
Здесь $D_s(\omega)d\omega$ --- полное число мод данной ветви
в интервале частот. При переходе к частоте множитель
$dk/d\omega=1/c_s$ необходимо учесть дополнительно к $k^{d-1}$.

\subsection*{Низкотемпературная энергия}

Пусть $k_{\mathrm B}T$ мало по сравнению с энергией на границе
линейной области. Верхний предел интеграла можно продолжить до бесконечности:
\begin{align*}
    U_{\mathrm{th}}
       &=\sum_s\int_0^\infty D_s(\omega)
          \frac{\hbar\omega}{e^{\hbar\omega/(k_{\mathrm B}T)}-1}\,d\omega\\
       &=\frac{V_d S_{d-1}}{(2\pi)^d}
          \sum_s\frac{(k_{\mathrm B}T)^{d+1}}{\hbar^d c_s^d}
          \int_0^\infty\frac{x^d}{e^x-1}\,dx.
\end{align*}
Разложение $(e^x-1)^{-1}=\sum_{n=1}^\infty e^{-nx}$ даёт
\[
    \int_0^\infty\frac{x^d}{e^x-1}\,dx
       =\sum_{n=1}^\infty\frac{\Gamma(d+1)}{n^{d+1}}
       =\Gamma(d+1)\zeta(d+1).
\]
Следовательно,
\[
    C_V=\frac{V_d S_{d-1}}{(2\pi)^d}(d+1)\Gamma(d+1)\zeta(d+1)
       k_{\mathrm B}\left(\frac{k_{\mathrm B}T}{\hbar}\right)^d
       \sum_s c_s^{-d}.
\]
\rresult{Для линейного акустического спектра $C_V\propto T^d$.}

\subsection*{Явные результаты}

Для одной ветви одномерной цепочки ($V_1=L$, $S_0=2$)
\[
    \frac{U_{\mathrm{th}}}{L}
       =\frac{\pi}{6}\frac{(k_{\mathrm B}T)^2}{\hbar c},
    \qquad
    \frac{C_V}{L}=\frac{\pi}{3}
       \frac{k_{\mathrm B}^2T}{\hbar c}.
\]
Для двумерной системы с линейными ветвями ($V_2=A$, $S_1=2\pi$)
\[
    \frac{U_{\mathrm{th}}}{A}
       =\frac{\zeta(3)}{\pi}
          \frac{(k_{\mathrm B}T)^3}{\hbar^2}\sum_s\frac1{c_s^2},
    \qquad
    \frac{C_V}{A}=\frac{3\zeta(3)}{\pi}
          \frac{k_{\mathrm B}^3T^2}{\hbar^2}\sum_s\frac1{c_s^2}.
\]
Для трёхмерного изотропного кристалла с одной продольной и двумя
поперечными ветвями ($V_3=V$, $S_2=4\pi$)
\begin{align*}
    \frac{U_{\mathrm{th}}}{V}
       &=\frac{\pi^2}{30}\frac{(k_{\mathrm B}T)^4}{\hbar^3}
          \left(\frac1{c_l^3}+\frac2{c_t^3}\right),\\
    \frac{C_V}{V}
       &=\frac{2\pi^2}{15}\frac{k_{\mathrm B}^4T^3}{\hbar^3}
          \left(\frac1{c_l^3}+\frac2{c_t^3}\right).
\end{align*}
Степень температуры определяется размерностью, а коэффициент
зависит от числа поляризаций и скоростей звука.

\subsection*{Дебаевское обрезание и высокие температуры}

В модели Дебая заменяем зону Бриллюэна сферой того же объёма.
Для каждой ветви одноатомной решётки должно остаться $N$ мод:
\[
    \frac{V_d S_{d-1}}{d(2\pi)^d}k_D^d=N,
    \qquad \omega_{D,s}=c_sk_D,
    \qquad \Theta_{D,s}=\frac{\hbar\omega_{D,s}}{k_{\mathrm B}}.
\]
С учётом конечного обрезания
\[
    C_V=\frac{V_d S_{d-1}}{(2\pi)^d}
       k_{\mathrm B}\sum_s
       \left(\frac{k_{\mathrm B}T}{\hbar c_s}\right)^d
       \int_0^{\Theta_{D,s}/T}
          \frac{x^{d+1}e^x}{(e^x-1)^2}\,dx.
\]
При $T\gg\max_s\Theta_{D,s}$ каждая мода даёт $k_{\mathrm B}$,
так что $C_V\to dNk_{\mathrm B}$ для $d$ ветвей идеальной
$d$-мерной модели. Для трёхмерного кристалла это
\rconcept{закон Дюлонга--Пти}: $C_V\to3Nk_{\mathrm B}$.
Если движения атомов не ограничены плоскостью, двумерная геометрия
не уменьшает число степеней свободы с трёх до двух.

При одинаковой дебаевской частоте в трёхмерной модели
\[
    C_V=9Nk_{\mathrm B}\left(\frac T{\Theta_D}\right)^3
       \int_0^{\Theta_D/T}\frac{x^4e^x}{(e^x-1)^2}\,dx,
\]
откуда
\[
    C_V\simeq\frac{12\pi^4}{5}Nk_{\mathrm B}
       \left(\frac T{\Theta_D}\right)^3
       \quad(T\ll\Theta_D).
\]

\answer
Акустический вклад при низких температурах подчиняется законам
$C_V\propto T$, $T^2$, $T^3$ в размерностях $1$, $2$, $3$ соответственно.
При классическом возбуждении всех мод теплоёмкость равна
$k_{\mathrm B}$, умноженной на их число.
\endgroup

\clearpage
\task[2]{Вклад оптических фононов}

\condition
Определить вклад в теплоёмкость от оптических фононов.
Пренебречь дисперсией оптической ветви: $\omega_{\boldsymbol k}=\omega_0$.
Пусть $\mathcal N_{\mathrm{opt}}$ --- полное число таких нормальных мод.

\solution

\subsection*{Модель Эйнштейна}

Каждая мода является осциллятором одной и той же частоты. Обозначим
$\Theta_E=\hbar\omega_0/k_{\mathrm B}$. Тогда
\begin{align*}
    U_{\mathrm{opt}}&=\mathcal N_{\mathrm{opt}}\hbar\omega_0
       \left(\frac12+\frac1{e^{\Theta_E/T}-1}\right),\\
    C_{\mathrm{opt}}&=\mathcal N_{\mathrm{opt}}k_{\mathrm B}
       \left(\frac{\Theta_E}{T}\right)^2
       \frac{e^{\Theta_E/T}}{(e^{\Theta_E/T}-1)^2}.
\end{align*}
При низкой температуре оптический фонон требует конечной энергии,
поэтому число возбуждений экспоненциально мало:
\[
    C_{\mathrm{opt}}\simeq\mathcal N_{\mathrm{opt}}k_{\mathrm B}
       \left(\frac{\Theta_E}{T}\right)^2e^{-\Theta_E/T},
       \qquad T\ll\Theta_E.
\]
При $T\gg\Theta_E$ восстанавливается классический предел
\[
    C_{\mathrm{opt}}\longrightarrow\mathcal N_{\mathrm{opt}}k_{\mathrm B}.
\]

\subsection*{Подсчёт мод и область применимости}

В трёхмерном кристалле с $r$ атомами в ячейке и $N_c$ ячейками
имеется $3r$ ветвей: три акустические и $3r-3$ оптических.
Если всем оптическим ветвям приписать одну частоту, то
$\mathcal N_{\mathrm{opt}}=(3r-3)N_c$.
При высоких температурах сумма обоих вкладов равна
$3rN_ck_{\mathrm B}$.

\rnote{Экспоненциальное подавление связано с щелью, а не с
полным отсутствием дисперсии.} Для диспергирующей оптической ветви
с минимумом $\omega_{\min}>0$ сохраняется множитель
$e^{-\hbar\omega_{\min}/(k_{\mathrm B}T)}$,
но степенной предэкспоненциальный множитель зависит от формы минимума.

\answer
В приближении Эйнштейна оптический вклад определяется формулой выше;
при $T\ll\Theta_E$ он экспоненциально мал, при $T\gg\Theta_E$
каждая оптическая мода даёт $k_{\mathrm B}$.

\Needspace{16\baselineskip}
\task[3]{Изгибная ветвь графена}

\condition
Найти вклад в теплоёмкость от одной нижней фононной ветви
двумерного листа графена с дисперсией $\omega=\beta k^2$, $\beta>0$.
Обсудить температурную зависимость.

\solution

\subsection*{Плотность состояний квадратичной ветви}

Для одной ветви в плоскости
\[
    d\mathcal N=\frac A{(2\pi)^2}2\pi k\,dk.
\]
Поскольку $d\omega=2\beta k\,dk$, получаем
\[
    D(\omega)=\frac A{4\pi\beta}.
\]
Плотность состояний постоянна. Она не пропорциональна $\omega$,
как у линейной двумерной ветви.

\subsection*{Энергия и теплоёмкость}

При температуре много меньше характерной энергии верхней границы ветви
\begin{align*}
    U_{\mathrm{ZA,th}}
       &=\frac A{4\pi\beta}\int_0^\infty
          \frac{\hbar\omega\,d\omega}{e^{\hbar\omega/(k_{\mathrm B}T)}-1}\\
       &=\frac A{4\pi\beta}\frac{(k_{\mathrm B}T)^2}{\hbar}
          \int_0^\infty\frac{x\,dx}{e^x-1}
        =\frac{A\pi}{24}\frac{(k_{\mathrm B}T)^2}{\hbar\beta}.
\end{align*}
Следовательно,
\[
    \frac{C_{\mathrm{ZA}}}{A}
       =\frac\pi{12}\frac{k_{\mathrm B}^2T}{\hbar\beta}.
\]
\rresult{Квадратичная изгибная ветвь даёт линейную теплоёмкость.}
Линейные продольная и поперечная ветви в плоскости дают $T^2$,
поэтому в этой модели изгибная ветвь доминирует при достаточно малых $T$.

\subsection*{Область применимости}

Результат относится к \rassumption{гармонической модели свободного ненатянутого листа}.
Натяжение может сделать спектр линейным при самых малых $k$;
подложка и ангармонизм также меняют низкочастотный спектр.
В конечном образце континуальный закон требует
$\hbar\beta(2\pi/L)^2\ll k_{\mathrm B}T$.

\answer
Для одной ветви $\omega=\beta k^2$
\[
    U_{\mathrm{ZA,th}}=\frac{A\pi k_{\mathrm B}^2T^2}{24\hbar\beta},
    \qquad C_{\mathrm{ZA}}=\frac{A\pi k_{\mathrm B}^2T}{12\hbar\beta}.
\]

\end{document}
